%%%%%%%%%%%%%%%%%%%%%%%%%%%%%%%%%%%%%%%%%%%%%%%%%%%%%%%%%%%%%%%%%%%%%%
% Lớp 4 | Chủ đề 1. Ôn tập và bổ sung | Giải bài toán có ba bước tính
% Dạng 4D1?5-2: Tính ngược từ cuối
%%%%%%%%%%%%%%%%%%%%%%%%%%%%%%%%%%%%%%%%%%%%%%%%%%%%%%%%%%%%%%%%%%%%%%


% nguồn: 2026-11-18_lop4_chu-de-1-tu-luan (nhập 18/11/2026)
\begin{bt}%[4D1H5-2]%[ma:5ac2ea43]
% nguon: Toán tư duy lớp 4 – Bài 5 SGK Toán 4 KNTT
	Tìm một số, biết rằng lấy số đó cộng với $15$, được bao nhiêu đem nhân với $3$ thì được $120$.
	\loigiai{
		Tính ngược từ cuối, làm phép tính ngược lại (nhân thành chia, cộng thành trừ):\\
		Trước khi nhân với $3$: $120 : 3 = 40$.\\
		Trước khi cộng $15$: $40 - 15 = 25$.\\
		Số cần tìm là $25$. Thử lại: $(25 + 15) \times 3 = 120$.
	}
\end{bt}

% nguồn: 2026-11-18_lop4_chu-de-1-tu-luan (nhập 18/11/2026)
\begin{bt}%[4D1V5-2]%[ma:a6b1adb4]
% nguon: Toán tư duy lớp 4 – Bài 5 SGK Toán 4 KNTT
	Lan có một số kẹo. Lan cho em $5$ cái, rồi cho bạn một nửa số kẹo còn lại thì Lan còn $8$ cái. Hỏi lúc đầu Lan có bao nhiêu cái kẹo?
	\loigiai{
		Tính ngược từ cuối:\\
		Lan còn $8$ cái là một nửa số kẹo sau khi cho em, nên sau khi cho em Lan có $8 \times 2 = 16$ (cái).\\
		Lúc đầu Lan có: $16 + 5 = 21$ (cái).\\
		Đáp số: $21$ cái kẹo.
	}
\end{bt}

% nguồn: 2026-11-18_lop4_chu-de-1-tu-luan (nhập 18/11/2026)
\begin{bt}%[4D1C5-2]%[ma:3f38eb70]
% nguon: Toán tư duy lớp 4 – Bài 5 SGK Toán 4 KNTT
	Một người bán trứng. Lần thứ nhất bán một nửa số trứng và thêm $1$ quả. Lần thứ hai bán một nửa số trứng còn lại và thêm $1$ quả thì vừa hết. Hỏi lúc đầu người đó có bao nhiêu quả trứng?
	\loigiai{
		Tính ngược từ lần bán cuối.\\
		Lần hai: bán một nửa và thêm $1$ quả thì hết, nghĩa là $1$ quả đó chính là nửa còn lại. Trước lần hai có $1 \times 2 = 2$ (quả).\\
		Lần một: sau khi bán một nửa và thêm $1$ quả thì còn $2$ quả, nên nửa số trứng lúc đầu là $2 + 1 = 3$ (quả).\\
		Lúc đầu có $3 \times 2 = 6$ (quả).\\
		Thử lại: bán $3 + 1 = 4$ quả, còn $2$; bán $1 + 1 = 2$ quả, hết.
	}
\end{bt}

% nguồn: 2026-11-18_lop4_chu-de-1-tu-luan-bo-sung (nhập 18/11/2026)
\begin{bt}%[4D1H5-2]%[ma:4245cf77]
% nguon: Toán tư duy lớp 4 – Bài 5 SGK Toán 4 KNTT
	Tìm một số, biết rằng lấy số đó nhân với $4$ rồi trừ đi $36$ thì được $100$.
	\loigiai{
		Tính ngược từ cuối (trừ thành cộng, nhân thành chia):\\
		Trước khi trừ $36$: $100 + 36 = 136$.\\
		Trước khi nhân $4$: $136 : 4 = 34$.\\
		Số cần tìm là $34$. Thử lại: $34 \times 4 - 36 = 136 - 36 = 100$.
	}
\end{bt}
